\selectlanguage{ngerman}%
\frenchspacing%
\begin{ledgroupsized}[r]{120mm}%
\footnotesize%
\pstart%
\noindent\textbf{Überlieferung:}%
\pend%
\end{ledgroupsized}%
\begin{ledgroupsized}[r]{114mm}%
\footnotesize%
\pstart%
\parindent -6mm%
\makebox[6mm][l]{\textit{L}}%
Aufzeichnung:
LH XXXV~1,\;9 Bl.\;30.
Ein Zettel;
Ränder zum Teil beschnitten und unregelmäßig.
Eineinhalb Seiten;
Textfolge:
Bl.\;\href{\Scan{LH_35_01_09_030v}}{30\,v\textsuperscript{o}},%
~\href{\Scan{LH_35_01_09_030r}}{30\,r\textsuperscript{o}}.
\pend%
\end{ledgroupsized}%
\vspace{5mm}%
\begin{ledgroup}%
\footnotesize%
\pstart%
\noindent%
\textbf{Datierungsgründe:}
Die Datierung der Aufzeichnung N.\;\ref{39978} beruht auf inhaltlichen Betrachtungen.
Siehe die ausführliche Begründung auf S.\;\pageref{39978_dat}.%
\pend%
\end{ledgroup}%
\selectlanguage{latin}%
\frenchspacing%
\vspace{8mm}%
\normalsize%
\pstart%
\noindent%
\bracks{%
\href{\Scan{LH_35_01_09_030v}}%
{30\,v\textsuperscript{o}}%
}
\edlabel{LH_35_01_09_030_030v_vdf-1}%
Inter principia demonstrandi%
\indexS{principium demonstrandi}
in Geometria%
\indexS{geometria}
et Mechanica%
\indexS{mechanica}
\edtext{%
perutilia,%
\indexS{principium perutile in geometria et mechanica}
etiam%
}{%
\lemma{perutilia,}%
\Bfootnote{%
\textit{(1)}~et %
\textit{(2)}~etiam%
~\textit{L}%
}}
hoc est ex natura continui%
\indexS{natura continui}%
\indexS{continuum}
a me sumtum,%
\indexS{principium sumtum e natura continui}
quod in rebus continue variabilibus,%
\indexS{res variabilis continue}
paulatim non momento
quantitas aliqua acquiritur
vel amittitur.%
\edlabel{LH_35_01_09_030_030v_vdf-2}%
\indexS{quantitas acquisita vel amissa paulatim}
\pend%
\begin{edfigure}[1.750em][1.50em]%
\includegraphics[scale=1.05]%
{gesamttex/edit_VIII-4/images/LH_35_01_09_030_d_030v.pdf}
\edcaption[0.5em]{%
\bracks{\textit{Fig.\:1}}%
}%
\end{edfigure}
\pstart%
Sic demonstro
\edlabel{LH_35_01_09_030_030v_prdem-1}%
\edtext{%
(contra Libellatorum metallicorum praxin)%
\indexS{praxis libellatorum metallicorum}%
\indexS{libellator metallicus}%
}{%
\lemma{(contra}%
\Bfootnote{%
\textit{(1)}~Libellatores %
\textit{(2)}~Libellatorum %
\bars{metallicorum \textit{erg.}} praxin)%
~\textit{L}%
}%
\lemma{Libellatorum metallicorum praxin}%
\Cfootnote{%
Im wissenschaftlich-technischen Sprachgebrauch der Frühneuzeit waren \textit{libellatores} im Allgemeinen die "Vermesser".
Im Bergbau bezeichnete dieser Begriff die Markscheider.
}%
}%
\edlabel{LH_35_01_09_030_030v_prdem-2}
funiculum \textit{CD}
ad horizontem inclinatum%
\indexS{funiculum inclinatum}%
\indexS{horizon}%
\bracks{,}
a
\edtext{%
pondere in medio~\textit{B} appenso%
\indexS{pondus appensum funiculo}%
}{%
\lemma{pondere}%
\Bfootnote{%
\textit{(1)}~\textit{BC} %
\textit{(2)}~in medio \textit{B} appenso%
~\textit{L}%
}}
inaequaliter tendi,%
\indexS{funiculum tensum inaequaliter}
magis supra in \textit{BD} %
quam infra in \textit{BC}.
Quoniam
si filum \textit{DC} continue erigi ponatur,%
\indexS{filum inclinatum continue erectum}
donec ex inclinato fiat%
\indexS{filum factum ex inclinato verticale}
\edtext{%
vertical\bracks{e},%
\indexS{filum verticale}%
}{%
\lemma{verticali}%
\Bfootnote{%
\textit{L\;ändert Hrsg.}%
}}
utique in transitu illo ultimo%
\indexS{transitus ultimus}
ubi perpendiculare fit,%
\indexS{filum perpendiculare}
sola pars \textit{BD} tenditur
a
\edtext{%
pondere \textit{B},%
\indexS{pondus appensum funiculo}%
}{%
\lemma{pondere}%
\Bfootnote{%
\textit{(1)}~\textit{BC} %
\textit{(2)}~\textit{B},%
~\textit{L}%
}}
pars autem inferior \textit{BC}
nullo modo.
Quod
si ergo antea aequalis fuisset utriusque tensio%
\indexS{tensio funiculi inclinati}
sub quacunque inclinatione%
\indexS{inclinatio funiculi tensi}%
\bracks{,}
necesse
\edtext{%
esset
momento transitus%
\indexS{momentum transitus}
totam medietatem tensionis%
\indexS{medietas tensionis}
ab uno amitti,
ab altero aequari.%
}{%
\lemma{esset}%
\Bfootnote{%
\textit{(1)}~\angles{vel} %
\textit{(2)}~momento transitus %
\bracks{...} altero aequari.%
~\textit{L}%
}}
\bracks{%
\href{\Scan{LH_35_01_09_030r}}%
{30\,r\textsuperscript{o}}%
}
\pend%
\newpage%
\pstart%
Quod
si quis asserat
aliquandiu aequalem manere tensionem%
\indexS{tensio funiculi inclinati}
sub aliquibus inclinationibus,%
\indexS{inclinatio funiculi tensi}
inde alicubi incipere decrescere,
v.\,g.
sub inclinatione 45 graduum?%
\indexS{inclinatio 45 graduum}
Ejus initii ratio reddi non potest.%
\indexS{ratio initii reddenda}
Cur citra illum limitem%
\indexS{limes}%
\bracks{,}
nullum fiat
\edtext{%
tensionis%
}{%
\lemma{tensionis}%
\indexS{tensio funiculi inclinati}%
\Bfootnote{%
\textit{erg.\;L}%
}}
decrementum%
\indexS{decrementum tensionis}
inter mutandum
\edtext{%
situm;%
\indexS{situs mutandus}%
}{%
\lemma{situm}%
\Bfootnote{%
\textit{erg.\;L}%
}}
supra illum fiat?
\pend%
\vspace{2.0em}%
\pstart%
\noindent%
\centering%
*     *     *
\pend%
\vspace{2.0em}%
\selectlanguage{german}%
\frenchspacing%
\footnotesize%
\pstart%
\noindent%
\textls{Ausführliche Datierungsgründe:}%
\pagelabel{39978_dat}
Der Träger der vorliegenden Aufzeichnung N.\;\ref{39978} weist weder Spuren eines Wasserzeichens noch irgendein äußeres datierungsrelevantes Merkmal auf.
Das "aus der Natur des Kontinuums" gewonnene und "in der Geometrie und der Mechanik überaus nützliche" Prinzip, das im Mittelpunkt der Aufzeichnung steht, ist zweifelsohne das Kontinuitätsprinzip (vgl. oben, \refpassage{LH_35_01_09_030_030v_vdf-1}{LH_35_01_09_030_030v_vdf-2}).
Dessen heuristische Anwendung als \textit{principium demonstrandi}, wie sie in der Aufzeichnung N.\;\ref{39978} exemplifiziert wird, lässt sich im Leibniz-Nachlass bis zum Anfang 1678 zurückverfolgen;
vgl. etwa \textit{De corporum concursu}, scheda~III, lemma~1 (\textit{LSB} VIII,\;3 N. 58\textsubscript{4}, S.\;563\,f., bes. S.\;563.19--564.3).\cite{01401}
Es ist auszuschließen, dass die vorliegende Aufzeichnung N.\;\ref{39978} vor Anfang 1678 entstand.
Denn anders als in \textit{LSB} VIII,\;3 N. 58\textsubscript{4}\cite{01401} wird das Kontinuitätsprinzip in der Aufzeichnung N.\;\ref{39978} als selbständiger und bereits "angenommener" (\textit{sumptum}) Grundsatz dargestellt.
\pend%
\pstart%
Die erste öffentliche Ankündigung des Kontinuitätsprinzips erfolgte wohl in dem \textit{Extrait d’une lettre sur un principe general utile à l'explication des loix de la nature par la consideration de la sagesse divine}, den Leibniz im Juliheft 1687 der \textit{Nouvelles de la république des lettres} drucken ließ (S.\;744–753; \textit{LPG}~III, S.\;51–55).\cite{02054}
Auch der von Leibniz in den letzten Monaten des Jahres 1688 in Wien\indexO{Wien} verfasste Entwurf \textit{Principium quoddam generale non in mathematicis tantum sed et physicis utile} (\textit{LSB} VI,\;4 N.\;371),\cite{02055} thematisiert das Kontinuitätsprinzip ausführlich.
In beiden Texten wird der heuristische Nutz dieses Prinzips in den strengen Wissenschaften -- in den mathematischen sowie in den physischen -- ausdrücklich hervorgehoben.
Inhaltlich stimmt die vorliegende Aufzeichnung N.\;\ref{39978} mit diesem Ansatz überein und könnte zu einer ähnlichen oder späteren Zeit entstanden sein.
\pend%
\pstart%
Diese Feststellung lässt sich ebenfalls durch die Anspielung auf die \textit{Libellatorum metallicorum praxis} im Text (\refpassage{LH_35_01_09_030_030v_prdem-1}{LH_35_01_09_030_030v_prdem-2}) bestätigen.
Im Staatsarchiv Hannover\indexO{Hannover} waren bis 1943 zwei Manuskripte von Leibnizens Hand erhalten, welche die Überschriften \textit{Chordae tensae libellator}\bracks{\textit{um}} und \textit{Libellatio} trugen und somit wohl eben jener "Tätigkeit der Markscheider" gewidmet waren, die in der Aufzeichnung N.\;\ref{39978} erwähnt wird;
vgl. NLA~HA~MS Nr.\;59 (früher: Hannover, HStA Ms.\;E\;3) Bl.\;330 u.\;331.\cite{02553}\cite{02554}
Beide Manuskripte, die heute als Kriegsverlust gelten,
waren auf den 29.\;August (8.\;September) 1686 datiert.
Es ist zu vermuten,
dass die Aufzeichnung N.\;\ref{39978} im Zusammenhang mit diesen Texten verfasst wurde, oder zumindest nicht vor ihnen.
Weitere heute verschollene Texte des Hannoveraner\indexO{Hannover} Staatsarchivs zeigten, dass Leibniz sich im gleichen Sommer -- wohl an der Hand von Nikolaus \textsc{Voigtel}s gerade veröffentlichter \textit{Geometria subterranea} (Eisleben 1686)\cite{02556} -- intensiv mit der Markscheidekunst und der mit ihr verwandten Wägekunst beschäftig haben soll;
siehe hierzu \textsc{R.\;Koser}, "Über die Sammlung von Leibniz-Handschriften im Staatsarchiv zu Hannover",
\ledpb%
\textit{Sitzungsberichte der königlich preußischen Akademie der Wissenschaften zu Berlin} XXV (1902), S.\;565.\cite{02555}
Demgemäß wäre eine Anspielung auf die Tätigkeit der \textit{libellatores metallici} wie die, die in der Aufzeichnung N.\;\ref{39978} (\refpassage{LH_35_01_09_030_030v_prdem-1}{LH_35_01_09_030_030v_prdem-2}) vorkommt, vor dem Sommer 1686 weniger wahrscheinlich gewesen.
Hieraus ergibt sich der wahrscheinlichste Terminus post quem der Datierung.%
\pend%
\pstart%
Ein plausibler Terminus ante quem lässt sich aus folgender terminologischer Beobachtung gewinnen.
Das Kontinuitätsprinzip wird in der Aufzeichnung N.\;\ref{39978} zwar unmissverständlich angesprochen, aber noch nicht mithilfe eines spezifischen "Fachbegriffs".
Auch angesichts des bescheidenen Umfangs der Aufzeichnung hätte Leibniz einen solchen Begriff wohl verwendet, wenn dieser sich in seinem Sprachgebrauch bereits etabliert hätte.
In dem nur handschriftlich überlieferten zweiten Teil seines \textit{Specimen dynamicum} etwa bedient er sich durchaus des "Fachbegriffs" \textit{lex continuitatis}, wenn er auf das Kontinuitätsprinzip Bezug nimmt (vgl. etwa \textit{LGM}~VI, S.\;249).\cite{02557}\cite{01043}
Dieser Umstand kann als Zeichen gedeutet werden, dass die Aufzeichnung N.\;\ref{39978} zu einem früheren Zeitpunkt entstand.
Den zweiten Teil von Leibnizens \textit{Specimen dynamicum} hatte Otto Mencke\indexN{\textls{Mencke} (Menken, Menkenius, Menque), Otto 1644\textendash1707} für eine Veröffentlichung im Maiheft 1695 der \textit{Acta eruditorum} vorgesehen,\cite{01023} wie dies sich seinem Brief an Leibniz vom 13./23.\;März desselben Jahres entnehmen lässt (\textit{LSB} I,\;11 N.\;238 S.\;342.6--8).\cite{02558}
Zu dieser Veröffentlichung ist es freilich nie gekommen:
Mitte September und Anfang November 1695 wartete Mencke\indexN{\textls{Mencke} (Menken, Menkenius, Menque), Otto 1644\textendash1707} noch vergeblich darauf, dass Leibniz den zweiten Teil des \textit{Specimen dynamicum} übersenden würde (vgl. \textit{LSB} I,\;11 N.\;471, S.\;695.6--11;\cite{02559} I,\;12 N.\;109, S.\;126.9--12).\cite{02560}
Unter der Voraussetzung aber, dass der zweite Teil des \textit{Specimen dynamicum} ungefähr zur gleichen Zeit entstand wie der im Aprilheft 1695 der \textit{Acta eruditorum} veröffentlichte erste Teil (S.\;145--157;\cite{02032} \textit{LMG}~VI, S.\;234--246)\cite{01043} oder nicht wesentlich später, dürfte auch die Aufzeichnung N.\;\ref{39978} spätestens im Frühjahr 1695 verfasst worden sein, wahrscheinlich aber früher.
\pend%
\pstart%
Aus den angegebenen Gründen ergibt sich schließlich für die vorliegende Aufzeichnung N.\;\ref{39978} folgende Gesamtdatierung:
Eine Entstehung vor September 1686 (jedenfalls nach 1677) ist nicht auszuschließen, aber unwahrscheinlich;
N.\;\ref{39978} dürfte auch nicht viel später verfasst worden sein, jedenfalls nicht nach dem Frühjahr 1695; eine spätere Entstehung ist aber möglich.
\pend%
\selectlanguage{latin}%
\frenchspacing%
\normalsize%
